\section{The exterior comparison and the full signed objective}
\label{sec:assembly}
We prove the exterior calculation of \cite{assembly}, with its radial
premise supplied by Theorem~\ref{thm:radial}, using the direct competitive
argument \cite{direct}. Let $f$ be any measurable assignment satisfying
$S_{-r}(f)\le S_{-r}(g)$, and put $\rho=-r$. Use the notation of Analytic interface
\ref{prior:localization}, and write $R=\calR(B)$ and $d=d(B)$.
If $d\le1/m$, \eqref{eq:local-noise} proves the desired conclusion.
Throughout this section assume
\begin{equation}\label{eq:exterior}
                              d>1/m.
\end{equation}

\subsection{The normalized cost and its exterior floor}
From $0\le\delta\le2D$, one has
$r\abs{\Delta}=\abs{D-\delta}\le D$; therefore
$W/w=1+O(r/L)$, including when $W>w$.
Since $f$ is admissible for $F_p(B)$, $F_p(B)\ge W$.
The balanced optimal offsets are a dual test at $p$, so
\[
 R\le\frac{w-W}{2}-h_U\cdot(p-U)
       \le\frac\Delta2+CL\sqrt{R\delta}.
\]
Here \eqref{eq:offset-input} is scaled using
\eqref{eq:radial-identity}, and the entropy bound gives the required
fixed row cap. Young's inequality implies, for large $m$,
\begin{equation}\label{eq:Rpaid}
              0\le R\le\Delta+CL^2\delta
                   \le\Delta+\delta/r=D/r=:J.
\end{equation}
The division by $r$ is indispensable.

Let $q=d/\alpha$ and $e=\norm{BB^T/\alpha^2-P}_{\Frob}$.
If $s_i$ are the singular values of $B/\alpha$ on $\one^\perp$,
including zeros, Procrustes minimization gives
$q^2=\sum_i(s_i-1)^2$ and $e^2=\sum_i(s_i^2-1)^2$; hence $e\ge q$.
Also $c/\sqrt L\le q\le C\sqrt m$ by \eqref{eq:exterior} and
$W\asymp w$. Equation~\eqref{eq:global-covariance-floor} gives
\begin{equation}\label{eq:exteriorfloor}
 J\ge R\ge\frac{cq}{m^3}\ge\frac{c}{m^3\sqrt L},
                   \qquad d^2\le C m^{3/2}L R.
\end{equation}
These estimates concern the score remainder, not a replacement of the
actual assignment by an affine one.

Put $s=\sqrt{W/w}$ and $\widehat B=B/s$. By
Corollary~\ref{cor:unnormalized},
$E_{\rm rad}(\widehat B)\le CL^2R\le CrL^2/m\le CL^2/m^2$.
The radial identity gives $\calR(\widehat B)\le CR$, and
\eqref{eq:one-row} yields
\[
 \max_i\left|\frac{\abs{\widehat B_i}^2}{\beta^2}-1\right|
           \le C\sqrt{\frac{m^2R}{L}}
           \le C\sqrt{\frac{mr}{L}}=o(1).
\]
The universal operator bound \eqref{eq:actual-operator} scales to
$\norm{\widehat B}_{\rm op}^2\le C/m$. Thus every hypothesis of
Lemma~\ref{lem:proxy} is proved for this arbitrary competitive assignment.
It gives an exactly centered equal-row matrix $Z$, with median gradient
at most $C\sqrt r$.
Since $\norm{B-\widehat B}^2=w(s-1)^2\le2R$,
\begin{equation}\label{eq:proxyerror}
              E_Z:=\norm{Z-B}^2\le CL^2R\le CL^2J.
\end{equation}

\subsection{A maximum comparison with an explicit absolute tail}
Write $Y_i=Z_i\cdot X/\beta$, $M=\max_iY_i$, and let $M_*$ be the
regular unit-row maximum. Centering gives $M\ge0$ pointwise.
Set $a_0=(m-1)/m$ and let
$\nu=\sqrt{a_0}M+G/\sqrt m$, with $G$ independent standard Gaussian.
The underlying Gaussian score covariance is
$a_0C+\one\one^T/m$, where $C=ZZ^T/\beta^2$.
It has unit diagonal and dominates $\one\one^T/m$. The regular version
has covariance $I_m$. Analytic interface~\ref{prior:max} therefore gives
\[
 H(t):=\Prob\{\nu\le t\}-\Prob\{\max_iZ_i^{\rm iid}\le t\}\ge0.
\]
Set $\varepsilon_{\rm free}=\E M_*-\E M\ge0$. Integration of the CDFs
and cancellation of the independent common Gaussian yield
\[
 \int H=\sqrt{a_0}\varepsilon_{\rm free},\qquad
 \E M_*^2-\E M^2=\frac2{a_0}\int tH(t)\dd t.
\]
For $t\ge T$, $H(t)\le m\overline\Phi(t)$; for $t\le-T$,
$H(t)\le\Phi(t)$. Splitting the integral and using $H\ge0$ gives,
for $T\ge1$,
\begin{equation}\label{eq:maxsecond}
       (\E M_*^2-\E M^2)_+
           \le CT\varepsilon_{\rm free}+Cm e^{-T^2/2}.
\end{equation}
No deficit-relative tail estimate is assumed.
With $V=\beta M=\max_i Z_i\cdot X$ and $V_*=\beta M_*$,
$\E V_*=w$ and \eqref{eq:maxsecond} becomes
\begin{equation}\label{eq:scaledmax}
 (\E V_*^2-\E V^2)_+
       \le C\beta T(w-\E V)+C\beta^2m e^{-T^2/2}.
\end{equation}
In particular $\E V\le w$.

\subsection{The actual assignment pays its test-score error}
Write $z_i(X)=Z_i\cdot X$ and
$c_f(X)=V-\sum_i f_i z_i\ge0$. Both $\E c_f$ and $w-\E V$ are
nonnegative, and their sum is exactly
\begin{equation}\label{eq:routingproxy}
 w-\ip ZB=\frac{\Delta+E_Z}{2}\le CL^2J.
\end{equation}
Here and below $\ip ZB=\sum_iZ_i\cdot B_i$.
This identity is valid even if $\Delta<0$ and the masses are unequal.
Pointwise, $z_i\le V$ and $V\ge0$, so
\[
 V^2-\sum_i f_i z_i^2
      =\sum_i f_i(V-z_i)(V+z_i)\le2Vc_f.
\]
If $A_X=\max_i\abs{Y_i}$, then $c_f\le2\beta A_X$ and
$V\le\beta A_X$. Unit marginal Gaussian tails imply
\[
 \E[A_X^2\one_{\{A_X>T\}}]
                  \le Cm(T^2+1)e^{-T^2/2}.
\]
Splitting according to $V\le\beta T$ gives
\begin{equation}\label{eq:assignmentsquare}
 \E V^2-\E\sum_i f_i z_i^2
       \le2\beta T\E c_f+C\beta^2m(T^2+1)e^{-T^2/2}.
\end{equation}
Take $T=\sqrt{32L}$. Combining
\eqref{eq:scaledmax}--\eqref{eq:assignmentsquare} and
$\beta T\le CL/m$ gives
\begin{equation}\label{eq:score-square}
 \E\sum_i f_i z_i^2
       \ge\E V_*^2-\frac{CL^3J}{m}-C\beta^2mL e^{-16L}.
\end{equation}

\subsection{A quadratic certificate with unequal actual masses}
Define
\[
 Q_i=\E[f_i(XX^T-I)],\qquad
 T_i=Z_iZ_i^T-\frac{Z^TZ}{m}.
\]
Since $\sum_iQ_i=0$, every test row has squared norm $\beta^2$, and
$\sum_i p_i=1$, the numerator is
\begin{equation}\label{eq:tensor-numerator}
 T_f:=\sum_i\ip{Q_i}{T_i}_{\Frob}
       =\E\sum_i f_i z_i^2-\sum_i p_i\abs{Z_i}^2
       =\E\sum_i f_i z_i^2-\beta^2.
\end{equation}
Thus exact balance is not required. The denominator satisfies
\begin{equation}\label{eq:tensor-denominator}
 \sum_i\norm{T_i}_{\Frob}^2
 =m\beta^4-\frac{\tr((ZZ^T)^2)}m
 \le D_*:=\frac{\alpha^4(m-1)(m-2)}m,
\end{equation}
because $\tr(ZZ^T)=w$ and $\rank(ZZ^T)\le m-1$.
Consequently $A_2(f)=\frac12\sum_i\norm{Q_i}_{\Frob}^2
\ge(T_f)_+^2/(2D_*)$ by Cauchy's inequality. This remains true if
the actual test tensor vanishes.

For the regular reference,
\[
 T_*:=\E V_*^2-\beta^2
      =\alpha^2\{\E[(\max_i Z_i^{\rm iid})^2]-1\}>0,
 \qquad A_2(g)=\frac{T_*^2}{2D_*},\qquad T_*/D_*\le Cm.
\]
For any $e\ge0$, $(T_*-e)_+^2\ge T_*^2-2T_*e$; hence
\eqref{eq:score-square}--\eqref{eq:tensor-denominator} imply
\[
 A_2(f)-A_2(g)\ge-CL^3J-CL^2e^{-16L}.
\]
The exterior floor \eqref{eq:exteriorfloor} pays the absolute error:
$L^2e^{-16L}/J\le C m^{-13}L^{5/2}=o(1)$. Therefore
\begin{equation}\label{eq:exteriorA2}
                         A_2(f)-A_2(g)\ge-CL^3J.
\end{equation}
This is an exterior estimate for the localized assignment. It is not a
global relative second-chaos inequality near the regular partition.

\subsection{The entire signed higher tail}
Define $\calH_r(f)=\sum_{k\ge3}(-r)^kA_k(f)$, and let $g_0$ be the
regular assignment associated with a closest regular centroid $B_0$.
Its linear routing loss is
\[
 R_0=w-\ip{B_0}B=\frac{\Delta+d^2}{2}\ge0.
\]
Let $\varepsilon=\frac12\norm{f-g_0}_{L^1}$.
On each point the total mass sent away from the regular winning label
pays at least the top-two gap. Equation~\eqref{eq:regular-margin},
followed by the elementary bathtub estimate, gives
\begin{equation}\label{eq:epsilon-routing}
                  \varepsilon\le\min\{1,C\sqrt{mR_0}\}.
\end{equation}
This holds for fractional assignments too.

Put $\tau^2=1/(4L)$. Hypercontractivity and $p_i\le3/(2m)$ give
\[
 \norm{T_\tau f}_2^2+\norm{T_\tau g_0}_2^2\le C/m.
\]
For the difference, write $e_i=\E\abs{f_i-g_{0,i}}\le C/m$.
Since $\abs{f_i-g_{0,i}}\le1$,
\[
 \norm{T_\tau(f_i-g_{0,i})}_2^2
    \le e_i^{2/(1+\tau^2)}\le Ce_i/m.
\]
Summation yields $\norm{T_\tau(f-g_0)}_2^2\le C\varepsilon/m$.
On degrees at least three, factor the exact signed multiplier through
$T_\tau$; the residual operator has norm at most $(r/\tau^2)^3$.
Cauchy's inequality therefore gives
\begin{equation}\label{eq:HC-tail}
       \abs{\calH_r(f)-\calH_r(g_0)}
                         \le Cr^3L^3\frac{\sqrt\varepsilon}{m}.
\end{equation}
All higher degrees remain in this estimate, with their signs; no finite
chaos truncation is substituted.

If $J\ge d^2$, then $R_0\le J$, and
\eqref{eq:epsilon-routing}--\eqref{eq:HC-tail} and the exterior floor give
\[
 \frac{\abs{\calH_r(f)-\calH_r(g_0)}}J
 \le\frac{Cr^3L^3}{m^{3/4}J^{3/4}}
 \le Cr^3m^{3/2}L^{27/8}.
\]
If $J<d^2$, then $R_0\le d^2$. Use $d=\alpha q$, $J\ge cq/m^3$,
and $q\ge c/\sqrt L$ to obtain
\[
 \frac{\abs{\calH_r(f)-\calH_r(g_0)}}J
 \le Cr^3m^{7/4}L^{13/4}q^{-1/2}
 \le Cr^3m^{7/4}L^{7/2}.
\]
The latter also dominates the first case. Rotation invariance of the
regular Hermite weights gives
\begin{equation}\label{eq:tail-final}
 \abs{\calH_r(f)-\calH_r(g)}
                       \le Cr^3m^{7/4}L^{7/2}J.
\end{equation}

\subsection{Completion of the proof}
Hermite orthogonality gives exactly
\[
 S_{-r}(f)-S_{-r}(g)
       =D+r^2\{A_2(f)-A_2(g)\}+\calH_r(f)-\calH_r(g).
\]
Equations~\eqref{eq:exteriorA2} and \eqref{eq:tail-final}, with $J=D/r$,
therefore imply
\begin{equation}\label{eq:strict-exterior}
 S_{-r}(f)-S_{-r}(g)
   \ge D\{1-CrL^3-Cr^2m^{7/4}L^{7/2}\}\ge D/2>0
\end{equation}
for sufficiently large $m$, uniformly in $0<r\le1/(m-1)$.
The two endpoint errors are $O(L^3/m)$ and
$O(L^{7/2}/m^{1/4})$. Strictness follows from $D=rJ>0$ in the exterior.
This contradicts competitiveness. Consequently every competitive assignment
lies inside $d\le1/m$, where Analytic interface~\ref{prior:local-global} forces
the regular assignment. Noncompetitive assignments already have strictly
larger noise value. This proves both the inequality and its equality
classification in Theorem~\ref{thm:main}, directly for all measurable
simplex-valued assignments. No full-noise attainment, purity, stationarity, interface
regularity, or limiting optimization argument is required. All thresholds
are independent of finite $N$ and of positive $r$ in the voting interval.
\qed
